\section{Introduction}
\label{sec:intro}

Among the convex polytopes with $n$ vertices inscribed in the unit sphere
$\mathbb{S}^{d-1}$, which has the largest volume? Fejes T\'{o}th posed the
three-dimensional case in \emph{Regular Figures}, where an inequality already
settles $n = 4, 6, 12$ as the regular tetrahedron, octahedron and icosahedron
\cite[p.~263]{fejes-toth-1964}, as restated in \cite{horvath-2016-survey}. The
remaining cases have resisted a uniform treatment, and progress has come case
by case.

In three dimensions the line runs as follows. Goldberg stated a conjecture on
the shape of the extremal polyhedron for the dual isoperimetric problem
\cite{goldberg-1935}; Grace transferred the question to the inscribed setting
and ran an early computer search \cite{grace-1963}; Bowen and Fisk supplied the
isomorph-free enumeration of sphere triangulations that made case analysis
possible \cite{bowen-fisk-1967}. Berman and Hanes then introduced the local
optimality condition now called Property Z, proved that a maximiser is
simplicial, and settled $n = 5, 7, 8$ \cite{berman-hanes-1970}. The next case
took fifty-six years: Hoehner and Ledford resolved $n = 9$, the maximiser being
a triaugmented triangular prism of volume $3\sqrt{2\sqrt{3}-3}$
\cite{hoehner-ledford-2026}. Beyond this the record is numerical. Mutoh
reported a computer search for $4 \leq n \leq 30$ \cite{mutoh-2003}, and
Hardin, Sloane and Smith tabulate putatively optimal configurations for
$4 \leq n \leq 130$ \cite{hardin-sloane-smith}; the values at $n = 10$ and
$n = 11$ are recorded as conjectures in \cite{oeis-a339262,oeis-a339263}. No
optimality proof is known for $n = 10$, $n = 11$ or any $n \geq 13$.

In higher dimensions Horv\'{a}th and L\'{a}ngi extended Property Z and
simpliciality to every dimension and solved $n = d+2$ outright, $n = d+3$ for
odd $d$, and $n = d+3$ for even $d$ among noncyclic polytopes; they also showed
that the maximal-volume inscribed cross-polytope is the regular one
\cite{horvath-langi-2016}. Exhaustive work in four dimensions additionally
needs to know which combinatorial types admit an inscribed realisation at all,
the problem Steiner raised in 1832 \cite{steiner-1832}. Hodgson, Rivin and
Smith characterised it in three dimensions
\cite{hodgson-rivin-smith-1992,dillencourt-smith-1996}, but no comparable
characterisation is known in higher dimensions \cite{padrol-ziegler-2016};
Firsching's inscribed rational realisations supply the answer for $d = 4$,
completely for $n \leq 8$ \cite{firsching-2017}.

This work presents an experimental methodology rather than new mathematics. A
single implementation, applied uniformly instead of case by case, reproduces
and verifies results spread across these papers, and every number it reports is
labelled with the strength of the evidence behind it. Claims about our own
computations and code carry a second kind of reference, written [R$k$] and
listed under \emph{Code and data} \artref{reproduce}. In summary:

\begin{enumerate}
  \item The published record is reproduced on $\mathbb{S}^{2}$ for
    $4 \leq n \leq 12$ and on $\mathbb{S}^{3}$ for
    $n \leq 8$ (Sections~\ref{sec:results3} and \ref{sec:results4}). At $d = 4$,
    $n = 7, 8$ the sweep ranges over every combinatorial type, where the
    published optimality results are restricted to a class, so it covers a
    wider class of polytopes, though only numerically.
  \item Every case carries an evidence level recomputed from its own per-type
    data rather than asserted in prose, so a claim cannot drift ahead of the
    numbers supporting it (Section~\ref{sec:method}).
  \item We note a failure mode affecting relaxation-based approaches. The
    signed-volume polynomial attains spurious maxima at self-intersecting
    configurations which satisfy Property Z exactly, and ascent converges to
    them from a large share of starts; on such a type every valid bound on the
    bare polynomial lies above the true maximum and so excludes nothing
    (Section~\ref{sec:trap}). The inequality that survives still supports a
    triage of combinatorial types.
  \item In four dimensions the sweep winners for $6 \leq n \leq 8$ are convex
    hulls of two regular polygons in orthogonal planes, so their volumes share
    one closed form (Section~\ref{sec:results4}).
\end{enumerate}

Nothing reported here is a certified bound, and the distinction between what is
computed and what is proved is recorded throughout.
